% !PRINTED 232
\weberpage{232}

%% [running head, not body text -- commented out 2026-08-07] Page furniture (running
%% head, folio, signature) is excluded from this edition by standing policy.
%% 232 \hfill Sechster Abschnitt. \hfill \S.\,74.

%<<CONT
wir annehmen, dass eine andere von den drei Grössen $B_{0,3}, B_{1,3}, B_{2,3}$
von Null verschieden ist.

\indent Um nun die Gleichungen (5), (6) in symmetrischer Weise zu
behandeln und namentlich zu erkennen, welche Quadratwurzel
zu ihrer Lösung erfordert wird, verfahren wir so. Wir setzen
zur Abkürzung
\weq{7}{
\begin{array}{rcl}
B_0 &=& \alpha_0\,B_{0,0} + \alpha_1\,B_{0,1} + \alpha_2\,B_{0,2} \\
B_1 &=& \alpha_0\,B_{1,0} + \alpha_1\,B_{1,1} + \alpha_2\,B_{1,2} \\
B_2 &=& \alpha_0\,B_{2,0} + \alpha_1\,B_{2,1} + \alpha_2\,B_{2,2}\,.
\end{array}
}
Dann können wir die Gleichungen (5) und (6) so darstellen:
\weq{8}{
\begin{array}{rcccccl}
\alpha_0\,B_0 & + & \alpha_1\,B_1 & + & \alpha_2\,B_2 & = & 0 \\
\alpha_0\,B_{0,3} & + & \alpha_1\,B_{1,3} & + & \alpha_2\,B_{2,3} & = & 0\,.
\end{array}
}
Wir multipliciren die zweite mit einem unbestimmten Factor
$\alpha_3$ und addiren sie zur ersten, wodurch wir erhalten
\[
\text{(9)} \qquad \alpha_0\,(B_0 + \alpha_3\,B_{0,3}) + \alpha_1\,(B_1 + \alpha_3\,B_{1,3}) + \alpha_2\,(B_2 + \alpha_3\,B_{2,3}) = 0\,,
\]
eine Gleichung, die erfüllt ist, wenn wir setzen
\weq{10}{
\begin{array}{rcl}
B_0 + \alpha_3\,B_{0,3} &=& 0 \\
B_1 + \alpha_3\,B_{1,3} &=& \sigma\,\alpha_2 \\
B_2 + \alpha_3\,B_{2,3} &=& -\,\sigma\,\alpha_1\,.
\end{array}
}
Aus diesen drei Gleichungen, in Verbindung mit der Gleichung
(6) haben wir die Unbekannten $\alpha_0, \alpha_1, \alpha_2, \alpha_3, \sigma$ zu bestimmen.
Wir setzen die Gleichungen zunächst ausführlich hierher
\weq{11}{
\begin{array}{rcrcrcrcl}
\alpha_0\,B_{0,0} & + & \alpha_1\,B_{0,1} & + & \alpha_2\,B_{0,2} & + & \alpha_3\,B_{0,3} & = & 0 \\
\alpha_0\,B_{1,0} & + & \alpha_1\,B_{1,1} & + & \alpha_2\,(B_{1,2} - \sigma) & + & \alpha_3\,B_{1,3} & = & 0 \\
\alpha_0\,B_{2,0} & + & \alpha_1\,(B_{2,1} + \sigma) & + & \alpha_2\,B_{2,2} & + & \alpha_3\,B_{2,3} & = & 0 \\
\alpha_0\,B_{3,0} & + & \alpha_1\,B_{3,1} & + & \alpha_2\,B_{3,2} & & & = & 0\,,
\end{array}
}
und wenn wir hieraus $\alpha_0, \alpha_1, \alpha_2, \alpha_3$ eliminiren, so ergiebt sich
eine quadratische Gleichung für $\sigma$, nach deren Lösung man die
Verhältnisse $\alpha_0 : \alpha_1 : \alpha_2 : \alpha_3$ aus linearen Gleichungen bestimmen
kann.

\indent Die quadratische Gleichung für $\sigma$ aber lautet in Determinantenform
\[
\begin{vmatrix}
B_{0,0}, & B_{0,1}, & B_{0,2}, & B_{0,3} \\
B_{1,0}, & B_{1,1}, & B_{1,2} - \sigma, & B_{1,3} \\
B_{2,0}, & B_{2,1} + \sigma, & B_{2,2}, & B_{2,3} \\
B_{3,0}, & B_{3,1}, & B_{3,2}, & 0
\end{vmatrix} = 0\,.
\]
